% (C) 2026 Ralf Stephan, in collaboration with Claude Code.
% Released under CC0 1.0 Universal (public-domain dedication).
% See https://creativecommons.org/publicdomain/zero/1.0/

\documentclass[11pt]{article}
\usepackage[T1]{fontenc}
\usepackage[utf8]{inputenc}
\usepackage{lmodern}
\usepackage[margin=2.5cm]{geometry}
\usepackage{amsmath,amssymb}
\usepackage{braket}
\usepackage[hidelinks]{hyperref}

% Lean names and paths: verbatim, breakable at . _ /
\DeclareUrlCommand\lean{\urlstyle{tt}}
\def\UrlBreaks{\do\.\do\_\do\/}
% Extra stretch for paragraphs with long, unbreakable Lean names
\emergencystretch=6em

\title{Open problems:\\ Projective measurements are equilibration processes}
\author{}
\date{}

\begin{document}
\maketitle

\noindent
Open problems raised by the technical brief in this directory, and by its Lean formalization
(see \lean{formalization-log.tex}).

\section*{From the brief}
\begin{enumerate}
\item \textbf{Unbounded observables.} For unbounded $X$ the Lindblad equation of EQ-LIND holds
  only on a domain; the brief claims that the blockwise solution (5) ``extends the resulting
  semigroup to every density operator'' and that $\rho(t)$ stays a density operator. Neither the
  extension nor Theorem~1 for unbounded $X$ is proved. (Brief, Sec.~3, after eq.~(5).)
\item \textbf{Non-commutative Scheff\'e lemma.} ``Hilbert--Schmidt convergence to a density
  operator implies trace-norm convergence'' for density operators is cited without reference or
  proof. The formalization bypasses it, so the lemma itself remains unproved here. (Brief,
  Sec.~3, between eqs.~(9) and~(10).)
\item \textbf{Continuous spectrum.} The claim that all four conditions fail on a continuous part
  of the spectrum of $X$: an operator commuting with $X$ and supported on a non-atomic spectral
  sector ``cannot be a non-zero trace-class operator'', and Lindblad fixed points there are not
  density operators. Not proved. (Brief, Sec.~4, ``Observables with continuous spectra''.)
\item \textbf{A formalism with finite entropy.} When $S(P_X(\rho)) = +\infty$, EQ-KL does not
  determine a closest ensemble. The brief suggests that ``the real problem is the Hilbert space
  formalism, as it does not guarantee finite entropy for all density operators''; a state space
  in which every ensemble has finite entropy is not given. (Brief, Sec.~4, ``Relative entropy in
  infinite dimensions''.)
\item \textbf{Other dynamics with the same limit.} ``Other dynamics may have the same asymptotic
  map'': which Lindblad (or other) dynamics have $\mathcal{P}_X$ as their long-time limit, and
  whether physical measurements are equilibrations ``that happen at a fast scale'', is left open.
  (Brief, Sec.~4, ``Equilibration as a black-box description''; Sec.~5.)
\item \textbf{Repeatability and state update.} The brief claims that repeatability is
  ``intrinsic in requiring that all outputs are equilibria of the process'', and that state update
  on a known outcome ``is equivalent to what happens in the classical case''. Neither is made
  precise or proved. (Brief, Sec.~5.)
\item \textbf{Destructive measurements.} ``The black-box idea could be extended to destructive
  measurements'', with the system state as input and ``the final state of the measurement device''
  as output. No such characterization is given. (Brief, Sec.~5.)
\end{enumerate}

\section*{From the formalization}
\noindent Beyond finite dimension the formalization has a part in any dimension, which covers
bounded observables with pure point spectrum only; it leaves open the items ``Unbounded
observables'' and ``Continuous spectrum'' above and ``Bridge from the spectral theorem'' below.
\begin{enumerate}
\item \textbf{Bridge from the spectral theorem.} Observables are given by an eigendecomposition
  $(P_x, x)$. That a self-adjoint $X$ with pure point spectrum has one, with $P_x = \mu_S(\{x\})$
  its spectral projections, is not formalized. (Log, ``Hilbert spaces of any dimension'',
  ``Differences from the brief''.)
\item \textbf{Basis independence of the relative entropy.} Lindblad's formula for $D(\rho\|\sigma)$
  is shown to be independent of the chosen eigenbases only when $S(\rho) < \infty$; for
  $S(\rho) = \infty$ this is open. (Log, ``Hilbert spaces of any dimension'', ``Differences from
  the brief''.)
\end{enumerate}

\end{document}
